\section{Medología}
Para garantizar coherencia, calidad y un buen trabajo en paralelo, se utilizó:

\begin{enumerate}[label=\textbf{\arabic*.}]
    \item \textbf{Formato de código estandarizado}
    
    \begin{itemize}
        \item Plantilla llamada \texttt{template.v} para tener una misma estructura.
        \item Uso de la herramienta \texttt{Verible} para formateo automático y consistencia de estilo en todo el código.
        \item Señales llamadas igual al estándar y en minúscula.
        \item Nombres de estados en mayúscula.
    \end{itemize}
    
    \item \textbf{Flujo de trabajo con GitHub}
    
    \begin{itemize}
        \item Desarrollo por \textit{branches} para funcionalidades aisladas.
        \item Manejo de tareas mediante \textit{issues}.
        \item Revisiones e integraciones de código mediante \textit{pull requests}.
    \end{itemize}
    
    \item \textbf{Repartición de módulos}
    
    \begin{itemize}
        \item Daniel Sáenz: Módulo Transmisor.
        \item Rodrigo Sánchez: Módulo Sincronizador.
        \item Brandon Jiménez: Módulo Receptor.
    \end{itemize}
\end{enumerate}